\chapter*{Resumen}
\addcontentsline{toc}{chapter}{Resumen}

Los agentes de programación basados en modelos de lenguaje pueden modificar un
repositorio, ejecutar pruebas y reaccionar a sus resultados. Cuando el objetivo
atraviesa varias capas, sin embargo, un único agente concentra demasiado
contexto y una división excesiva introduce coordinación, contratos y conflictos
de integración. Esta tesis estudia esa tensión, denominada \emph{paradoja de la
granularidad}, mediante el diseño, la implementación y la evaluación de
\sistema, un orquestador local de agentes de código.

El aporte técnico combina dos responsabilidades. Un planificador semántico
propone unidades de trabajo justificadas por el objetivo y el repositorio. Una
política determinista y versionada evalúa si conviene conservar cada unidad como
hoja cohesiva o aceptar un corte semántico propuesto. El resultado se compila en
un grafo híbrido con jerarquía de integración y relaciones tipadas para
artefactos, interfaces y conflictos. Cada intento se ejecuta en un
\emph{worktree} aislado; el orquestador inspecciona el diff, controla el alcance,
valida el commit candidato exacto, integra de abajo hacia arriba y publica sólo
un resultado acompañado por manifiesto y comprobante. La historia canónica del
run es un journal append-only de eventos de dominio.

La evaluación confirmatoria final se pre-registró como una serie pequeña y
acotada. Se utilizó un repositorio Node ESM con capas de dominio, aplicación y
API; dos formas de tarea ---una modificación multi-capa y una función cohesiva
de dominio---; dos repeticiones por forma; una única configuración adaptativa;
el modelo Codex \texttt{gpt-5.4-mini}; cero reintentos automáticos; y un oráculo
externo ejecutado sobre cada SHA candidato exacto. Un rehearsal previo verificó
el instrumento y quedó fuera del denominador.

Se fijaron dos hipótesis. \textbf{H-F1} exigía que las cuatro celdas llegaran a
\texttt{completed} y \texttt{delivered}, publicaran un SHA no vacío y pasaran el
oráculo externo. \textbf{H-F2} exigía que la tarea multi-capa produjera una raíz
compuesta con tres hojas ---dominio, aplicación y API---, y que la tarea cohesiva
produjera una sola hoja ejecutable. El resultado fue \textbf{PASS 4/4 para
H-F1} y \textbf{PASS 4/4 para H-F2}. Las cuatro entregas y topologías se derivan
de journals, snapshots, recibos y resultados de oráculo preservados junto con
el manuscrito.

El alcance de la conclusión es deliberadamente estrecho. La serie no tiene
controles A/B y usa un único target pequeño; por ello no demuestra superioridad
frente a no descomponer, causalidad, optimalidad, escalabilidad ni
generalización estadística. Las líneas históricas Warehouse, G5, G6, G7 y SP2
no se combinan con el resultado final: Warehouse permanece en 1/8, G5 produjo un
resultado comparativo negativo, G6 fue inconcluso, G7 no produjo candidatos y
SP2 fue piloto. Se conservan como evidencia adversa o formativa, no como piezas
incompletas de un denominador positivo.

La conclusión defendible es concreta: bajo la versión, el modelo, la política,
los límites y el protocolo congelados, \sistema\ realizó de forma reproducible
el recorrido de planificación, ejecución aislada, integración, validación
externa del commit exacto y entrega, y seleccionó las dos formas de topología
pre-registradas. El trabajo aporta además una disciplina de evidencia para no
confundir actividad de agentes, tests autocontenidos o resultados parciales con
software entregado y verificado.

\textbf{Palabras clave:} orquestación de agentes, modelos de lenguaje,
descomposición de tareas, granularidad adaptativa, worktrees de Git, validación
de commits, evidencia experimental.
